% ============================================
% Supplementary Note 2: Assembly Effect Analysis
% ============================================
\subsection*{Supplementary Note 2: Assembly Effect Analysis}
\label{sec:suppl:note2}

To understand the mechanistic basis of community-level selection (CLS), we compared coalescence outcomes to a ``direct assembly'' control. In coalescence, two parental communities are first assembled separately from the species pool, then mixed; in direct assembly, all species from both pools are mixed simultaneously without prior assembly history.

During the assembly process, species compete and some go extinct. Those that survive to coexistence tend to have weaker mutual interactions on average. We quantified this assembly effect by measuring the mean pairwise interaction strength $\langle \alpha_{ij} \rangle$ among species at different stages (Supplementary Fig. 26). Pre-assembly (red) represents the mean interaction strength among all species in the initial pool, before any competitive exclusion has occurred. Post-assembly (blue) represents the mean interaction strength among only those species that survived assembly to coexistence. Across all tested interaction strengths ($\mu = 0$ to $1.2$), post-assembly communities exhibited substantially reduced mean interaction strength compared to the pre-assembly pool. This reduction becomes more pronounced at higher $\mu$, where stronger competition drives more species to extinction, leaving only those with weaker mutual interactions.

The assembly process creates internal coherence within each parental community: species that ``survived together'' share a history of compatible (less competitive) interactions. We hypothesized that this shared history would increase the probability of CLS outcomes during coalescence, compared to direct assembly where no such history exists. To test this, we compared the distribution of coalescence outcomes between the two assembly methods. Simulation results (Supplementary Fig. 27) show that across five interaction strengths ($\mu = 0.2, 0.4, 0.6, 0.8, 1.0$), coalescence consistently produced higher fractions of CLS outcomes compared to direct assembly. At $\mu = 0.2$, coalescence yielded 37\% CLS vs direct 12\% CLS; at $\mu = 0.6$, coalescence yielded 59\% CLS vs direct 20\% CLS; at $\mu = 1.0$, coalescence yielded 65\% CLS vs direct 22\% CLS. Correspondingly, direct assembly showed much higher rates of Restructuring outcomes, where neither parental lineage dominates. This is expected because without assembly history, species from different pools have not been pre-filtered for compatibility.

Experimental results (Supplementary Fig. 28) corroborate these simulation findings. We applied the same analysis to experimental data across three nutrient conditions (Nutr$-$, Base, Nutr$+$). In the Base and Nutr$+$ conditions, coalescence showed elevated CLS rates compared to direct assembly, consistent with the simulation predictions. The Nutr$-$ condition showed more variable results, possibly due to the weaker competitive interactions in nutrient-limited environments.

These results support our proposed mechanism for CLS emergence. First, assembly acts as a filter, retaining species with compatible (weaker) interactions. Second, species within an assembled community share this compatibility, creating correlated retention during subsequent coalescence. Third, when two assembled communities coalesce, species from the same parent tend to survive or go extinct together, producing CLS outcomes. Fourth, without assembly history (direct assembly), no such correlations exist, leading to more Restructuring outcomes.

